\subsection{INTEGRAL FORMULA FOR THE REMAINDER IN TAYLOR'S FORMULA; PRIMITIVES OF HIGHER ORDER}

The formula for integration by parts of order n (II, p.~60, formula (11)) allows one to express in terms of an integral the remainder \(\mathbf{r}_{n}(x)\) in the Taylor expansion of order n of a function which admits a regulated \((n+1)^{th}\) derivative on an interval I (I, p.~22); indeed, on replacing, in (12), f by \(f'\) , b by x, and \(g(t)\) by the function \((t-x)^{n}/n!\) , it follows that

\[
\begin{array}{l} \mathbf {f} (x) = \mathbf {f} (a) + \mathbf {f} ^ {\prime} (a) \frac {(x - a)}{1 !} + \mathbf {f} ^ {\prime \prime} (a) \frac {(x - a) ^ {2}}{2 !} + \dots \\ \qquad + \mathbf {f} ^ {(n)} (a) \frac {(x - a) ^ {n}}{n !} + \int_ {a} ^ {x} \mathbf {f} ^ {(n + 1)} (t) \frac {(x - t) ^ {n}}{n !} d t \end{array}\tag{17}
\]

in other words

\[
\mathbf {r} _ {n} (x) = \int_ {a} ^ {x} \mathbf {f} ^ {(n + 1)} (t) \frac {(x - t) ^ {n}}{n !} d t,\tag{18}
\]

which formula often permits one to obtain simple bounds for the remainder.

Given a regulated function f on an interval I, an arbitrary primitive g of f, being continuous in I, admits a primitive in its turn; an arbitrary primitive of an arbitrary primitive of f is called a second primitive of f.~More generally, a primitive of a primitive of order n - 1 of f is termed a primitive of order n of f.~One sees immediately, by induction on n, that the difference of two primitives of order n of f is a polynomial of degree at most equal to n - 1 (with coefficients in E). A primitive of order n of f is entirely determined if one specifies its value and those of its first n - 1 derivatives at a point \(a \in I\) .

In particular, \(\int_{a}^{(n)}\mathbf{f}\) denotes that primitive of order \(n\) of \(\mathbf{f}\) which vanishes, together with its first \(n - 1\) derivatives, at the point \(a\) . The Taylor formula of order \(n - 1\) , applied to this primitive, shows that if \(\mathbf{g} = \int_{a}^{(n)}\mathbf{f}\) , then

\[
\mathbf {g} (x) = \int_ {a} ^ {x} \mathbf {f} (t) \frac {(x - t) ^ {n - 1}}{(n - 1) !} d t\tag{19}
\]

so reducing the determination of a primitive of order n to one single integral.

